\documentclass[11pt]{extarticle}
% \documentclass{article}[draft]
\usepackage[margin=1in]{geometry}
% \usepackage{style}

\usepackage{amsmath}
\usepackage{amsthm}
\usepackage{amssymb}
\usepackage{bbm}
\usepackage{xcolor}

% \input{macros.tex}
\newcommand{\comment}[1]{\textcolor{red}{(#1)}}

\begin{document}

\title{EC 2026 Rebuttal}
\author{}
\date{}
\maketitle

\newcommand{\kkcomment}[1]{\textcolor{magenta}{(#1 --KK)}}

\section{Concrete Questions Reponse Box}

\subsection*{R1Q1}

Thank you for your response.

\vspace{0.2cm}
\noindent 
R1Q1: You are correct that the IR property is defined ex-ante. We cannot hope for an 
ex-post IR guarantee because there are events on which an agent can have negative
% \comment{did you accidentally cut something off here}
utility due to competition. In our setting, \emph{ex-ante} IR simply states that it is beneficial for an agent to participate without knowing her or her competitor's value for each others' data. 
%\kkcomment{Explain further. Ex-post negative IR will raise red flags, so you want to provide further clarification.}

\vspace{0.2cm}
\noindent 
We avoided using the term ex-interim IR as we believed different 
interpretations could lead to confusion. We would first like to clarify what we believe
you mean by ex-interim IR. Traditionally the ex-interim stage is after an agent learns
their own type but before learning other agents' types. If we interpret types as 
value functions, then this definition does not make sense, as our setting focuses precisely on when agents do not know their values for each others' goods \emph{a priori}, but can learn them incrementally as they gradually receive the goods.

% ; this is because our problem
% focuses on when agents do not realize either their entire value function or their competitor's,
% but instead
% only portions of them depending on the quantity of goods exchanged.
This is distinct from the IC property, which states that it is never in either 
agent's interest to leave the algorithm early, even when they can use information acquired
from previous rounds to make this decision (see the factual corrections box).


% \comment{sep line}

% \vspace{0.2cm}
% \noindent 
% You are correct that the IR property is defined ex-ante. We cannot hope for an 
% ex-post IR guarantee because there are events on which an agent can have negative
% utility due to competition.

% \vspace{0.2cm}
% \noindent 
% This leaves considering ex-interim IR, which motivates our definition for IC. 
% An IC algorithm is defined by the following property: even when agents are allowed to 
% use information
% revealed from previous rounds, it is never in their interest to stop early.
% This is captured by the fact that each agent's stopping decisions are functions of
% the value random variables observed in previous rounds (see Definition 2.5).

% \vspace{0.2cm}
% \noindent 
% While ex-interim IR is typically defined as a non-negative conditional expected utility
% guarantee, the IC definition we give is equivalent.
% To see this, notice that IC implies that each agent's conditional expected 
% utility is non-negative
% (or else an agent could strategically exit early when this is violated and improve
% their expected utility, contradicting the inequality in Definition 2.5).
% Conversely, if an agents expected conditional utility is always non-negative there 
% is nothing to be gained by exiting early, which implies IC.

% I think you need to be more careful here. ex-interim is when you realize your (entire) valuation, but not the others'. So, an ex-interim definition does not make sense here since the whole point of the paper is to study when agents do not know their values. --KK

% \begin{itemize}

%     \item Yes IR is ex ante. In general can't have that its always IR because of the
%     randomness in data valuations.

%     \item
%     IC can be viewed as an ex interim guarantee
%     because of the stopping time forumulation. If you could leave early this 
%     would mean you could change your stopping time which by definition is a function
%     of the info you've seen.
%     the decision to stop early can be adaptive and in terms of the information seen! 
%     this is why its written in the language of changing to a different stopping time

% \end{itemize}  

\subsection*{R3Q1}

Thank you for your response.

\vspace{0.2cm}
\noindent 
R3Q1: We allow the random valuations and utilities to be very general to model the complex ways 
in which agents can value each other's goods and be in competition. For instance, 
we do not require that the value functions are strictly increasing. This is reflected in the real
world example of sharing data for machine learning where the value of data is 
tied to the accuracy of the model it produces: accuracy (and thus value), while generally trending 
upward as more data is obtained, does 
not necessarily increase strictly.
We assume that the value and utility functions are normalized to zero for convenience,
and boundedness is for technical reasons (and is a standard assumption).


% \comment{say practically not an issue}

\vspace{0.2cm}
\noindent 
We apologize for a poor phrasing of the technical condition in Theorem 3.5.
% In that condition, we
That condition was meant to
% say 
highlight
that the egalitarian welfare 
% satisfies the requirement 
implies our algorithm is IR
for every instance
of $(V^a,V^b,U^a,U^b)$, i.e. without any other assumptions, and even for pathological valuation functions (e.g. when valuations may decrease when you get more goods).
% \kkcomment{verify, and/or provide other pathological examples here.}).
The condition is also satisfied by the Nash welfare (and other popular welfares) in
non-pathological instances where maximizing welfare and both agents' utilities are not at odds. 


For instance, in
Example 1.1, replacing the Egalitarian welfare with the Nash welfare ends up yielding the same 
welfare maximizing algorithm: give agent $b$'s resource to agent $a$, and if agent $a$
observes value 1, give agent $b$ agent $a$'s resource. 
Moreover, the following calculations demonstrate that both agents have non-negative expected utility,
and thus the algorithm is IR. Following the notation in Example 1.1, let $s'$ denote the 
final exchange, then
\[
    \mathbb{E}\left[
        U^a_{s'}
    \right]
    =
    \frac{1}{2}(0)
    +
    \frac{1}{2}
    \left(
        \frac{1}{2}(3/4)
        +
        \frac{1}{2}(1)
    \right)
    =
    \frac{7}{16}
    \qquad
    \mathbb{E}\left[
        U^b_{s'}
    \right]
    =
    \frac{1}{2}(0)
    +
    \frac{1}{2}
    \left(
        \frac{1}{2}(3/4)
        +
        \frac{1}{2}(-1/4)
    \right)
    =
    \frac{1}{8} .
\]
We agree that this could be better explained in a revision.

% \vspace{0.2cm}
% \noindent 
% For the welfare function, on page 11 we highlighted the Egalitarian welfare, $W(x,y)=\min(x,y)$,
% because 
% $\mathbb{E}[\text{min}(X,Y)]\Rightarrow \mathbb{E}[X],\mathbb{E}[Y]\geq 0$
% and thus the implication in Theorem 3.5 is automatically satisfied. 
% However, the fact that $\mathbb{E}[XY]\not\Rightarrow \mathbb{E}[X],\mathbb{E}[Y]\geq 0$
% does not rule out the Nash welfare in general.
% While it is true that the Nash welfare may not satisfy the implication for every instance 
% of $(V^a,V^b,U^a,U^b)$, there are many instances for which it does. 
% For instance, in
% Example 1.1 replacing the Egalitarian welfare with the Nash welfare ends up yielding the same 
% welfare maximizing algorithm. 
% Moreover, under this algorithm, both agents' expected
% utilities are non-negative
% so our optimal algorithm is IR.
% The implication in Theorem 3.5 is meant as an auxiliary condition, to show that IR is 
% inevitable under some welfare functions, not as a restriction on the welfare functions 
% that can be used for a particular instance of $(V^a,V^b,U^a,U^b)$.
% given problem instance.

% both agents' expected utility 



% \vspace{0.2cm}
% \noindent 
% For instance, 

% \begin{itemize}
%     \item
%     utilities and valuations very general to model the complex ways in which agents 
%     can value each others' good and be in competition. 

%     \item We emphasized 
%     Egalitarian welfare $W(x,y)=\min(x,y)$ because satisfies 
%     $\mathbb{E}[\text{min}(X,Y)]\Rightarrow \mathbb{E}[X],\mathbb{E}[Y]\geq 0$
%     and thus the implication in Theorem 3.5 is automatically satisfied. 

%     However, while the Nash welfare may not satisfy this implication for every 
%     problem instance (check what I called this), there are setups where the nash welfare
%     still works.
%     For instance in example 1.1 the expected utilities of both agents are non negative
%     so the implication trivially holds.
% \end{itemize}  

\section{Factual Corrections Response Box}

\subsection*{Review 385A}

Review 385A:
``Even then, the definition still compares the expected utility of stopping early rather than 
an adaptive response to the information they have obtained up to this round.''

\vspace{0.2cm}
\noindent 
The definition of IC does allow for an adaptive response to the information an agent
has obtained up to the previous round.
In Definition 2.5, this is made explicit by the fact that any sequence of early stopping 
decisions
$\{\delta'_t\}_{t\geq 0}$ is comprised of functions which depend on the 
values realized by both agents from the previous rounds:
\(
    \{V^a_k\}_{k=0}^{i_t},
    \{V^b_k\}_{k=0}^{j_t}
\).
Therefore, the time an agent chooses to exit, defined as $T'$, 
depends on information received from
the previous round. 
% \kkcomment{can add here that it is information about her vaues and the others' values}.
% \comment{make clearer that the $V^a$ and $V^b$ are from 
% both users}
The requirement that 
\(
    \mathbb{E}\left[
        U^a_{s_{\min(T,T')}}
    \right]
    \leq 
    \mathbb{E}\left[
        U^a_{s_T}
    \right]
\)
says no matter which adaptive strategy is employed by agent $a$
to leave early, the agent is no better off than continuing to participate 
in the algorithm until it terminates at time $T$.

\subsection*{Review 385B}

Insert default response.

\subsection*{Review 385C}

Insert default response.

\section{Questions for Kirthi}

\begin{enumerate}
    % \item I think we can fold the factual inaccuracies for R1 and R3 into their 
    % question responses?

    % \item He kinda got us on the Nash welfare comment?

    \item For reviewer 1 should we quote the optional feedback as an innacuracy?
    Its the same point we're already addressing

    \item For reviewer 3 add a comment about how you could randomize who gets 
    the item first for the sake of fairness?
\end{enumerate}  

\end{document}
